\documentclass{beamer}

\mode<presentation>
{
  \usetheme{Singapore}      % or try Darmstadt, Madrid, Warsaw, ...
  \usecolortheme{crane} % or try albatross, beaver, crane, ...
  \usefonttheme{serif}  % or try serif, structurebold, ...
  \setbeamertemplate{navigation symbols}{}
  \setbeamertemplate{caption}[numbered]
} 

\usepackage[english]{babel}
\usepackage[autostyle=true]{csquotes}
\usepackage{xeCJK}
\usepackage[T1]{fontenc}
\usepackage{xcolor}
\usepackage{setspace}
\setmainfont{Charis SIL}
\setCJKmainfont{Noto Serif CJK HK}

\title[Pearl River Sinitic]{Sinitic Languages of the Pearl River Region \\ --- Unit 1.2}
\subtitle{13ᵗʰ EACL Summer School, Universität Stuttgart}
\author{Hilário de Sousa 蘇沙}
\institute{EHESS–CRLAO \\ {\tiny(membre associé)}}
\date{Monday 27 July 2026 \\ 丙午六月十四 \\ {\footnotesize \texttt{hilario@bambooradical.com}}}

\begin{document}

\begin{frame}
  \titlepage
\end{frame}

\section{Cantonese}

\begin{frame}[fragile]{Cantonese}

Cantonese < Canton \\
English \emph{\underline{C}anton} \textsuperscript{\emph{not canton}} < Portuguese \emph{Cantão} \\

\medskip
Português: \textsuperscript{\emph{since 16\textsuperscript{th} century}} \\
A província de Cantão 廣東 \\ 
A cidade de Cantão 廣州 \\
\smallskip
English: \\
\textcolor{gray!75}{The province of Canton 廣東
\textsuperscript{\emph{this use gradually faded away}}} \\
The \underline{city} of Canton 廣州

\end{frame}

\begin{frame}[fragile]{Cantonese}

\enquote{Cantonese} as an ethnic or linguistic term: \\
\enquote{Cantonese} \textsuperscript{\emph{usually}}≠ \enquote {related to the Province of 廣東}

\medskip
Sinitic dialect groups:
\begin{itemize}
    \item \large Yuè 粵
    \item Hakka 客家
    \item Mǐn 閩
    \item \small Unclassified Patois of N Guangdong 粵北土話
    \item \small Mandarin: Jūnhuà 軍話, SW Mandarin 西南官話
\end{itemize}
\small (+ 1\textasciitilde 1.5\%(?) of the population speak native non-Sinitic languages)

\end{frame}

\begin{frame}[fragile]{Cantonese}

\begin{figure}
\includegraphics[width=0.8\textwidth]{langatlaschinaguangdong}
\caption{\label{fig:langatlaschinaguangdong}Guangdong in the 1\textsuperscript{st} ed. of \emph{Language Atlas of China}}
\end{figure}

\end{frame}

\begin{frame}[fragile]{Cantonese}

\enquote{Cantonese} as linguistic or ethnic term: \\
\quad \enquote{Cantonese} \textsuperscript{\emph{usually}} only includes Yuè \\
\quad \enquote{Cantonese} \textsuperscript{\emph{usually}} excludes Hakka and Mǐn\\

\begin{figure}
\includegraphics[width=1\textwidth]{cantohakkateochew}
\end{figure}

\vfill{\tiny \url{https://www.youtube.com/watch?v=sYIW2rOnfss}}

\end{frame}

\begin{frame}[fragile]{Cantonese}

\enquote{Cantonese} as linguistic or ethnic term: \\
\quad \enquote{Cantonese} \textsuperscript{\emph{usually}}≠ \enquote {belonging to the Province of 廣東} \\
\quad \enquote{Cantonese} also spoken in large numbers in e.g. 廣西 \\ \quad \small \emph{(and most do not identify with 廣東)}\\

\begin{figure}
\includegraphics[width=1\textwidth]{guangxicanto}
\end{figure}
\vfill{\tiny \url{https://www.youtube.com/watch?v=B56nxfnSV4U}}

\end{frame}

\begin{frame}[fragile]{Cantonese}

Sinitic dialect groups of 廣西:
\begin{itemize}
    \item \large Yuè 粵
    \item \large SW Mandarin 西南官話 {\scriptsize (also Jūnhuà 軍話)}
    \item \normalsize Hakka 客家
    \item \normalsize Pínghuà 平話
    \item \small Xiāng 湘
    \item \small Mǐn (Southern) 閩南
\end{itemize}
\small (about 40\% of the population speak native non-Sinitic languages)

\end{frame}

\begin{frame}[fragile]{Cantonese}

\begin{figure}
\includegraphics[width=0.8\textwidth]{langatlaschinaguangxisinitic}
\caption{\label{fig:langatlaschinaguangxisinitic}Sinitic Guangxi in the 1\textsuperscript{st} ed. of \emph{Language Atlas of China}}
\end{figure}

\end{frame}

\begin{frame}[fragile]{Cantonese}

\enquote{Cantonese} includes what range of Yuè dialects?\\
\medskip
\texttt{\emph{Different people have different definitions}}\\
\medskip
Language Atlas of China: \enquote{Cantonese} = Yue dialect group\\

\medskip
\begin{quote}
    Four groups of Chinese dialects are spoken in Guangdong province.¹ They are \underline{Yue (Cantonese)}, Hakka (Kejia), Min and Shaozhou Tuhua.
\end{quote}
\medskip
\begin{quote}
    The Guangxi Zhuang Autonomous Region [...] the dialects spoken include the following six groups: Southwestern Mandarin, \underline{Yue (Cantonese)}, Hakka, Pinghua, Xiang and Min.
\end{quote}
{\tiny (the underlining is mine)}
\end{frame}

\begin{frame}[fragile]{Cantonese}

However, many Yuè dialects are not intelligible to the people of Canton \textsuperscript{\emph{without prior exposure}}.\\
e.g. YouTube Cantonese with Brittany {\tiny \url{https://www.youtube.com/@cantobritt}}:\\
{\scriptsize \emph{(Her understanding of Táishān 台山 Yuè is roughly on par with Hakka)}}

\begin{columns}
    \begin{column}{0.5\textwidth}
        \begin{figure}
            
\includegraphics[width=1\textwidth]{brittanyhoisan}
            {\tiny \url{https://www.youtube.com/watch?v=nuXg6KhpwXc}}
        \end{figure}    
    \end{column}
    \begin{column}{0.5\textwidth}
        \begin{figure}
            
\includegraphics[width=1\textwidth]{brittanyhakka}
            {\tiny \url{https://www.youtube.com/watch?v=z-PaOWHvWgQ}}
        \end{figure}    
    \end{column}
  \end{columns}


\end{frame}

\begin{frame}[fragile]{Cantonese}
De Sousa (2022), Cantonese:
\begin{itemize}
    \item Yuè dialect of Canton
    \item enclave Yuè varieties (\enquote{Cantonese enclaves}) resulted from the migration of people from the Canton area, primarily after the First Opium War (1840–1842)
    \item enclave Yuè varieties resulted from the migration of people from other Cantonese enclaves
\end{itemize}

\vfill{\tiny \setstretch{1} de Sousa, Hilário. 2022. The Expansion of Cantonese Over the Last Two Centuries. In Ye, Zhengdao (ed.), \emph{The Palgrave Handbook of Chinese Language Studies}, 409–440. Singapore: Springer Nature Singapore. \url{https://doi.org/10.1007/978-981-16-0924-4_35}\\}

\end{frame}


%\section{乜}

%\subsection{Tables and Figures}

%\begin{frame}{Tables and Figures}

%\begin{itemize}
%\item Use \texttt{tabular} for basic tables --- see Table~\ref{tab:widgets}, for example.
%\item You can upload a figure (JPEG, PNG or PDF) using the files menu. 
%\item To include it in your document, use the \texttt{includegraphics} command (see the comment below in the source code).
%\end{itemize}

%\begin{table}
%\centering
%\begin{tabular}{l|r}
%Item & Quantity \\\hline
%Widgets & 42 \\
%Gadgets & 13
%\end{tabular}
%\caption{\label{tab:widgets}An example table.}
%\end{table}

%\end{frame}

%\subsection{Mathematics}

%\begin{frame}{Readable Mathematics}

%Let $X_1, X_2, \ldots, X_n$ be a sequence of independent and identically distributed random variables with $\text{E}[X_i] = \mu$ and $\text{Var}[X_i] = \sigma^2 < \infty$, and let
%\[ S_n = \frac{X_1 + X_2 + \cdots + X_n}{n}
%      = \frac{1}{n}\sum_{i}^{n} X_i \]
%denote their mean. Then as $n$ approaches infinity, the random variables $\sqrt{n}(S_n - \mu)$ converge in distribution to a normal $\mathcal{N}(0, \sigma^2)$.

%\end{frame}

\end{document}
